\SourceLeaf{22}
\[
[1]\qquad \xi=\int\!\int\frac{\mathrm{d}p.\mathrm{d}q.\cos p.\sin p}{Ar}
\]
\par
Perinde obtinemus, per theorema sextum
\[
\begin{aligned}
[2]\qquad \xi&=-\int\!\int\frac{\mathrm{d}p.\mathrm{d}q.\sin p}{r^3}.(a-x)\bigl(\frac{(a-x)x}{AA}+\frac{(b-y)y}{BB}+\frac{(c-z)z}{CC}\bigr)
\end{aligned}
\]
\par
Denique theorema quartum nobis suppeditat
\[
\begin{aligned}
[3]\qquad &\int\!\int\frac{\mathrm{d}p.\mathrm{d}q.\sin p}{r^3}\bigl(\frac{(a-x)x}{AA}+\frac{(b-y)y}{BB}+\frac{(c-z)z}{CC}\bigr)=0\\
&\hspace{6em}\text{vel }=-\frac{4\pi}{ABC}
\end{aligned}
\]
prout punctum $M$ iacet vel extra corpus, vel intra corpus.

Iam quantitates $A,B,C$ tamquam valores particulares trium variabilium $\alpha,\beta,\gamma$ consideramus, ita comparatarum. ut $\alpha\alpha-\beta\beta$, $\alpha\alpha-\gamma\gamma$ sint constantes. Ita $\xi$ spectari poterit tamquam functio variabilium $\alpha,\beta,\gamma$ seu potius unius ex ipsis: variationes simultaneas quantitatum $\xi,\alpha,\beta,\gamma$ per characteristicam $\delta$ distinguemus. Facile concluditur ex aequatione $[1]$, crescentibus $\alpha,\beta,\gamma$ in infinitum, $\xi$ ultra omnes limites decrescere, quum manifesto vel valor minimus ipsius $r$ ultra omnes limites crescat. Statuere itaque oportet $\xi=0$ pro $\alpha=\infty$. Differentiando aequationem $[1]$ ita exhibitam
\[
\alpha\xi=\int\!\int\frac{\mathrm{d}p.\mathrm{d}q.\cos p.\sin p}{r}
\]
secundum characteristicam $\delta$, prodit
\[
\alpha\delta\xi+\xi\delta\alpha=-\int\!\int\frac{\mathrm{d}p.\mathrm{d}q.\cos p.\sin p.\delta r}{rr}
\]
Sed habemus
\[
\begin{aligned}
r\delta r&=-(a-x)\delta x-(b-y)\delta y-(c-z)\delta z\\
&=-(a-x)\cos p.\delta\alpha-(b-y)\sin p.\cos q.\delta\beta-(c-z)\sin p\sin q.\delta\gamma\\
&=-(a-x)x.\frac{\delta\alpha}{\alpha}-(b-y)y.\frac{\delta\beta}{\beta}-(c-z)z.\frac{\delta\gamma}{\gamma}\\
&=-\alpha\delta\alpha.\bigl(\frac{(a-x)x}{\alpha\alpha}+\frac{(b-y)y}{\beta\beta}+\frac{(c-z)z}{\gamma\gamma}\bigr)
\end{aligned}
\]
(propter $\alpha\delta\alpha-\beta\delta\beta=0$, $\alpha\delta\alpha-\gamma\delta\gamma=0$): hinc fit
\[
\alpha\delta\xi+\xi\delta\alpha=\delta\alpha.\int\!\int\frac{\mathrm{d}p.\mathrm{d}q.x\sin p}{r^3}\bigl(\frac{(a-x)x}{\alpha\alpha}+\frac{(b-y)y}{\beta\beta}+\frac{(c-z)z}{\gamma\gamma}\bigr)
\]
\par
Hinc subtrahendo aequationem $[2]$, in $\delta\alpha$ multiplicatam, postquam $A,B,C$ in $\alpha,\beta,\gamma$ mutatae sunt, fit
